\chapter{The Composite Soundtrack}\label{ch:composite-sound}

Everything so far has concerned light. Light can be divided in time, and glasses can select the parts. Sound cannot be divided that way. A loudspeaker fills the room, every ear in the room receives what it emits, and no shutter can admit one second of it and reject the next without making nonsense of both. A selective screening can show different people different pictures. Unless something more is done, it plays everyone the same sound.

This chapter treats sound as the harder multiplex. It sets out the four ways a variant family can handle sound and defends the one that was part of the idea from its beginning: a single score composed to serve every realization at once.

\section{Four arrangements}

\emph{Common sound} plays one soundtrack in the room for everyone. It keeps the room acoustically whole and constrains the pictures severely, because every realization must fit the same sound.

\emph{Individual sound} gives each viewer the soundtrack of their own realization through headphones, earpieces built into the glasses, or bone-conduction transducers. It allows the realizations to differ completely, and it turns the room into a set of private listening booths.

\emph{Hybrid sound} divides the soundtrack. Music and ambient effects play in the room; dialogue, or narration, or commentary, reaches each viewer individually. It keeps some of the room's shared sound and releases the realizations from sharing their words.

\emph{Complementary sound} plays one soundtrack in the room, like common sound, but composes it for the family rather than for any one realization. It is written to be right for every picture, meaning something different under each.

The first three are engineering choices. The fourth is an artistic one, and it is the one the original \emph{Protocols} pair proposed: two films watched at the same time with ``overlapping, complementary soundtracks.''

\section{What individual sound costs}

Individual sound is technically easy. Several audio channels can be broadcast to receivers in the glasses as readily as the timing signal, and the glasses remain devices that only output, recording nothing. The costs are not technical.

The first is isolation. A room in which everyone wears headphones is not quite a cinema. It is a set of people each hearing a private film, and much of what the Introduction called the shared occasion is gone. The picture is common, or at least commonly placed, but the sound, which carries more of a film's emotional timing than viewers usually notice, is not.

The second is subtler, and it applies to every arrangement: the audience is also a sound source. People laugh, gasp, go silent, shift in their seats. In an ordinary screening these sounds are part of the experience, and they are synchronized with it, because everyone is reacting to the same thing. In a selective screening they are not. The viewers of the comedy laugh at a moment when the viewers of the drama are watching a death. Headphones do not prevent this, because the laughter is in the room. The audience's own reactions are a shared channel that the film does not control.

The third cost is to the unfiltered view. A viewer without glasses sees the composite of the pictures. With individual sound, that viewer hears nothing of any realization's dialogue, or hears one realization's sound over everyone's pictures. The room's own version of the film, which \Cref{ch:unfiltered} argued could be designed, loses its soundtrack.

\section{The dialogue constraint}

Common and complementary sound both require the realizations to share whatever is audible, and the hardest thing to share is speech. Music can be ambiguous. Effects can serve several pictures. But words are words, and delivery is sound: the tone of a line is in the audio, not in the image.

This produces a precise constraint. If a face is seen speaking in any realization, the shared soundtrack must carry speech that matches its lips. So wherever two realizations both show a speaking face at the same moment, they must show it saying the same words, with the same timing. They can differ in what surrounds it, in framing, in reaction, in what the listener's face does, but not in the speech itself. Under shared sound, divergent dialogue is possible only where mouths are not visible: speech off screen, from behind, in wide shots, or through voice-over.

The constraint has a direct consequence for \Cref{ch:genre}. The second production strategy for a genre pair, the same blocking shot twice with different performances, relies on different delivery. Under shared sound, the delivery cannot differ. The pair must either rely on visual performance, expression and gesture and timing of reaction, while the lines themselves are spoken once, or move divergent dialogue off screen, or adopt hybrid sound. None of these is fatal, and the first is the most interesting. A line written so that one delivery can be heard as menace in one picture and as bluster in another is a small version of the complementary score, and good actors can give such readings.

\section{Complementary scoring}

The complementary score rests on a familiar fact about film music: its meaning depends on the picture. A slow, dissonant cue under a character climbing stairs in the dark is dread. The same cue under a character creeping to raid a refrigerator is comedy. Composers exploit this constantly, and audiences read music through what they see without noticing that they do. Studies of film sound have described the relation in both directions: music shapes how an image is read, often without being consciously heard \cite{gorbman}, and sound adds meaning to an image that viewers then attribute to the image itself, what Chion calls added value \cite{chion}. A complementary score relies on the same relation run once against several images.

A complementary score makes the exploitation systematic. It is written against several pictures at once, so that each cue is right under each realization. Formally, if $A_k$ is the set of soundtracks acceptable for realization $k$, the shared score must lie in their intersection,
\begin{equation}
S\in\bigcap_{k=1}^{m}A_k,
\end{equation}
and the art is in finding a score that is not merely acceptable under each picture but good under each. The intersection is small when the realizations want incompatible things from sound, loud where the other wants silence, fast where the other wants stillness, and it is larger than it looks when the realizations share the events that sound marks. The \emph{Protocols} pair is favorable: a thriller and a satire about the same regime want tension from the same moments, and a score of ominous, slightly overstated cues can be dread under one and parody under the other.

The complementary score does something neither individual nor hybrid sound can. It keeps the room acoustically single, so that everyone hears the same music at the same moment, and it makes that single sound mean different things to different people. It is the audio counterpart of the designed composite of \Cref{ch:unfiltered}, with one difference: the composite image is what the room sees only without glasses, while the complementary score is what everyone hears all the time. It is the most thoroughly shared element of a selective screening, and it is where the family is most visibly one work.

\section{Scheduling the audience}

The complementary score also suggests a remedy for the audience's own sounds. If the realizations are composed together, their emotional peaks can be placed together. A comic realization can put its laughs where the dramatic realization permits release, and a dramatic realization can place its silences where the comedy is between jokes. The envelopes of \Cref{ch:genre} give the structure for this: moments of shared tension and shared relief can be aligned at their boundaries, so that the room's laughter falls where nobody is being asked to grieve.

This cannot be made perfect, and it should not be. Some dissonance between what a viewer sees and what they hear the room doing is part of the experience of selective cinema. A viewer of the drama who hears laughter from part of the room at a solemn moment learns, audibly, that others are watching something else. That knowledge is a large part of what \Cref{ch:synchronized-difference} will call synchronized difference. The aim of scheduling is not to hide the difference but to keep it from becoming cruelty.

\section{Choosing an arrangement}

The four arrangements trade divergence for community. Individual sound permits the most divergence and keeps the least of the room. Common sound keeps the room and permits the least. Hybrid sound sits between them. Complementary sound keeps the whole room and permits more divergence than common sound, at the price of composing for the family from the start.

A family's sound arrangement should therefore follow from what kind of family it is. A tonal family, whose realizations share events, suits complementary sound well. A depth family, whose extra explanation is often verbal, pushes toward hybrid sound, or toward carrying explanation in the picture, as \Cref{ch:depth} noted. A family of endings can share its score until the branch and needs a composer who can end in several directions at once. A family of unrelated films on one screen, the weakest case in \Cref{ch:family}, needs individual sound and gives up most of what makes the form worth building.

The early notes behind this book did not argue for complementary sound. They simply assumed it: the two films would be watched simultaneously ``as they have overlapping, complementary soundtracks.'' The assumption turns out to carry the idea's social promise. A selective screening with shared, complementary sound is one event that people experience differently. A selective screening with private sound is several events that happen to share a room.
